% Part: second-order-logic
% Chapter: sol-and-sets
% Section: cardinalities

\documentclass[../../../include/open-logic-section]{subfiles}

\begin{document}

\olfileid{sol}{set}{crd}

\olsection{集合的势}

\begin{explain}
正如我们可以表达论域是有穷还是无穷、是可数还是不可数一样，我们也可以定义 $\Domain{M}$ 的子集具有有穷或无穷、可数或不可数的性质。
\end{explain}

\begin{prop}
公式 $\fn{Inf}(X) \ident$
\begin{multline*}
\lexists[u][(\lforall[x][\lforall[y][(\eq[u(x)][u(y)] \lif 
      \eq[x][y])]] \land {} \\\lexists[y][(X(y) \land \lforall[x][(X(x)
      \lif \eq/[y][u(x)])]])]
\end{multline*}
在变元指派~$s$ 下被满足，当且仅当 $s(X)$ 是无穷集。
\end{prop}

\begin{prop}
公式 $\fn{Count}(X) \ident $
\begin{multline*}
\lexists[z][\lexists[u][(X(z) \land 
    \lforall[x][(X(x) \lif X(u(x)))] \land {} \\ \lforall[Y][((Y(z) \land
      \lforall[x][(Y(x) \lif Y(u(x)))]) \lif X = Y])]])
\end{multline*}
在变元指派~$s$ 下被满足，当且仅当 $s(X)$ 是可数集。
\end{prop}

根据 Cantor 定理，我们知道存在不可数集，而且实际上，存在无穷多种不同层级的无穷大小。集合论发展出了一整套关于集合大小的算术，并将无穷基数分配给集合。自然数充当了衡量有穷集大小的基数。可数无穷集的势是第一个无穷势，称为~$\aleph_0$（“aleph-nought”或“aleph-zero”）。下一个无穷大小是~$\aleph_1$。它是一个集合在不为可数的（即大小为~$\aleph_0$）情况下所能具有的最小大小。我们可以将“$X$ 的大小为 $\aleph_0$”定义为
$\fn{Aleph}_0(X) \liff \fn{Inf}(X) \land \fn{Count}(X)$。$X$ 的大小为 $\aleph_1$，当且仅当它的所有子集都是有穷的或大小为~$\aleph_0$，但其本身大小不是~$\aleph_0$。因此，我们可以用公式 $\fn{Aleph_1}(X) \ident \lforall[Y][(Y \subseteq X \lif
  (\lnot\fn{Inf}(Y) \lor \fn{Aleph}_0(Y)))] \land \lnot
\fn{Aleph}_0(X)$ 来表达这一点。大小为 $\aleph_2$ 的定义方式类似，以此类推。

有一种大小特别令人关注，即所谓的连续统势。它是 $\Pow{\Nat}$ 的大小，或者等价地，是~$\Real$ 的大小。一个集合具有连续统势这一性质也可以在二阶逻辑中表达，但需要更多的工作。

\end{document}